\chapter{Sensitivitätsanalyse}\label{ch:ch3}
Der Hauptteil kann aus mehreren Kapiteln bestehen. Hier soll auf die Vorgehensweise sowie die Ergebnisse der Arbeit eingegangen werden. Dabei ist auf eine schlüssige und klare Strukturierung zu achten, die den roten Faden der Arbeit verdeutlicht. 

Kurzer einleitender Text nach jedem Kapitelbeginn und vor Beginn der \texttt{sections}. Es stehen also nie zwei Überschriften direkt hintereinander. An dieser Stelle soll dem Leser ein Überblick über das folgende Kapitel gegeben werden.

Hier habe ich etwas geändert
\section{Erstes Unterkapitel}\label{sec:ch3_foobar}
Dies wurde bereits von Fliess et~al. im Jahre 1995 gezeigt~\cite{Fliess.1995}.

Mit \texttt{biblatex}-Befehlen: Dies wurde bereits von \citeauthor{Fliess.1995} im Jahre~\citeyear{Fliess.1995} gezeigt~\cite{Fliess.1995}.


Das Logo des Instituts ist in Abbildung~\ref{fig:logo_isys1} dargestellt. 

  
\begin{figure} % avoid using additional specifiers for graphics position except for final tuning
 \centering    % preferable over \begin{center}...\end{center} as the latter includes space before and after the figure 
 \includegraphics[width=0.5\textwidth]{logos/logoISYS_deutsch}
 \caption{Put figure captions \emph{below} the figure.}
 \label{fig:logo_isys1} % always use \label AFTER \caption
\end{figure}

Die erzielten Ergebnisse sind in Tabelle~\ref{fig:example_table1} zusammengefasst.

\begin{table} % avoid using additional specifiers for graphics position except for final tuning
	\caption{Put table captions \emph{above} the table. Put units in header. Avoid vertical lines in a table.}
	\label{fig:example_table1}
	\centering
	\begin{tabular}{lll}
		\toprule
		Header & Header (\unit{unit}) & Header (\unitfrac{unit}{frac})   \\
		\midrule
		Content 1 & Content 2 & Content 3 \\
		Content 1 & Content 2 & Content 3 \\
		\bottomrule
	\end{tabular}
\end{table}

   

\section{Zweites Unterkapitel}


\subsection{Neuer Abschnitt}




